% Structural statistics, Minnesota k=4 test split, D_intra.
% Generated by evaluation/make_intra_table.py on 2026-08-14 from
%   results/samples/{B,C}_k4{,_nohop}.txt
% Requires: booktabs, float (for [H])

\begin{table}[H]
\centering
\caption{Structural statistics of generated supernode blocks against real
held-out blocks, on the Minnesota $k{=}4$ test split. Luxembourg is coarsened at
$k{=}7$ throughout; only the private graph is recoarsened, which brings the
private clusters onto the same size distribution as training (median 11 nodes in
both) and enlarges the shared test split from 47 supernodes to 114. For the upper
block, closer to the real column is better. Generation is one block per held-out
supernode, so all four configurations are scored on the same 114 blocks and the
real column is compared against itself, giving exactly $0$ on the lower block.
Resampling the real blocks to the same count gives a $W_1$ of
$0.0228 \pm 0.0108$ rather than $0$, which is the resolution limit of the metric
at this sample size.}\label{tab:dintra-k4-stats}
\begin{tabular}{lrrrrr}
\toprule
& \multicolumn{2}{c}{B: Lux only} & \multicolumn{2}{c}{C: Lux + Minn} & \\
\cmidrule(lr){2-3}\cmidrule(lr){4-5}
& hops & no hops & hops & no hops & Real \\
\midrule
Edges (median)                             & 10     & 12     & 11     & 11     & 11 \\
Edges (std)                                & 5.787  & 8.761  & 5.810  & 5.887  & 5.928 \\
Edges (mean)                               & 10.965 & 13.430 & 11.360 & 11.342 & 11.474 \\
Density (mean)                             & 0.2170 & 0.2359 & 0.2278 & 0.2206 & 0.2228 \\
Degree (median)                            & 2      & 2      & 2      & 2      & 2 \\
Degree (mean)                              & 1.8825 & 2.3057 & 1.9503 & 1.9473 & 1.9699 \\
Connected (\%)                             & 49.1   & 96.5   & 92.1   & 98.2   & 100.0 \\
Acyclic (\%)                               & 41.2   & 43.0   & 50.9   & 60.5   & 56.1 \\
\midrule
Degree distribution vs.\ real $\downarrow$ & 0.1054 & 0.3389 & 0.0527 & 0.0361 & 0.0000 \\
\bottomrule
\end{tabular}
\end{table}


% ---------------------------------------------------------------------------
% Companion table: paired bootstrap on the degree-distribution distance.
% ---------------------------------------------------------------------------

\begin{table}[H]
\centering
\caption{Paired bootstrap on the degree-distribution $W_1$, 2{,}000 resamples of
the 114 shared supernodes. Every configuration is scored on the same blocks in
the same order, so the supernode index is resampled and both sides recomputed
from one draw, removing the between-cluster variance. A negative difference
favours the right-hand configuration. Adding private training data is significant
in both condition variants; dropping the hop histogram is not, except in $B$,
where the comparison is confounded (see note).}\label{tab:dintra-k4-bootstrap}
\begin{tabular}{llrrr}
\toprule
Comparison & & $W_1$ left & $W_1$ right & 95\% CI on difference \\
\midrule
$B$ hops    & vs.\ $C$ hops    & 0.1054 & 0.0527 & $[-0.0943,\ -0.0064]$ \\
$B$ no hops & vs.\ $C$ no hops & 0.3389 & 0.0361 & $[-0.3743,\ -0.2051]$ \\
$C$ hops    & vs.\ $C$ no hops & 0.0527 & 0.0361 & $[-0.0473,\ +0.0160]$ \\
$B$ hops    & vs.\ $B$ no hops & 0.1054 & 0.3389 & $[+0.1362,\ +0.3174]$ \\
\bottomrule
\end{tabular}
\end{table}

% NOTE on the B hops vs B no hops row. The ablation's premise -- that the two
% models differ by h alone -- does not hold for B at k=4. The condition reports
% each graph's true locality, and B trains on Luxembourg only, so k has zero
% training variance and is left unstandardised; the Minnesota test records
% therefore enter at -3 where training was always 0. In the no-hop condition
% [n, rho_in, k] that is one of only three inputs, with no h columns to dilute
% it. C is unaffected: its training set contains both localities, so k is
% standardised normally. Read that row as h plus a k shift, not h alone.
